%%%%%
%
\section{S�lection de donn�es dans deux vecteurs selon des crit�res 
appliqu�s � l'un d'entre eux}

\itemstyle{Op\'erateur~:} \textsf{select\_inter\_data}
\index{select\_inter\_data}

%%%%%OBJET%%%%
\itemstyle{Objet~:}
On s�lectionne les donn�es de Vect$1$ et Vect$2$ si les donn�es de Vect$2$ sont 
dans un (voire deux) intervalles de validit� (c.� d. comprises 
entre crit$1$\_min et crit$1$\_max et �ventuellement entre crit$2$\_min et 
crit$2$\_max).

On peut ne renseigner qu'un seul crit�re, dans ce cas on passe pour les crit�res 
non renseign�s la valeur ($0.$/$0.$).

Les donn�es non s�lectionn�es sont mises � la valeur choisie par l'utilisateur
et pass�e en argument ("misval", $0$, autre valeur).

On obtient en sortie deux vecteurs de donn�es s�lectionn�es : $Vect1_{sel}$ et 
$Vect2_{sel}$. 

Cet op�rateur est bien adapt� � la s�lection de mesures de 
vent (DD,FF) sur des crit�res d'intervalle de validit� de la direction. 

%%%%%ARGUMENTS%%%%
\bigskip
%%%%%ENTREES%%%%
\begin{tabarg}
\multicolumn{3}{l}{\titretab {Entr\'ees}}\\
\hline
$Vect1$ & Vecteur, {\tt float} & Donn�es � s�lectionner (unit�
quelconque)\\
$Vect2$ & Vecteur, {\tt float} & Donn�es sur lesquelles porteront les tests (unit�
quelconque)\\
$crit1\_min$ & Coeff., {\tt float} & Premi�re borne min (m�me unit� que $Vect2$)\\
$crit1\_max$ & Coeff., {\tt float} & Premi�re borne max (m�me unit� que $Vect2$)\\
$crit2\_min$ & Coeff., {\tt float} & Deuxi�me borne min (m�me unit� que $Vect2$)\\
$crit2\_max$ & Coeff., {\tt float} & Deuxi�me borne max (m�me unit� que $Vect2$)\\
$val1$ & Coeff., {\tt float} & Valeur de remplacement des donn�es de Vect$1$ non
selectionn�es qui ne sont pas manquantes (m�me unit� que $Vect1$)\\
$val2$ & Coeff., {\tt float} & Valeur de remplacement des donn�es de Vect$2$ non
selectionn�es qui ne sont pas manquantes (m�me unit� que $Vect2$)\\
\end{tabarg}

\medskip
%%%%%SORTIES%%%%
\begin{tabarg}
\multicolumn{3}{l}{\titretab {Sorties}}\\
\hline
$Vect1_{sel}$ & Vecteur, \texttt{float} & Donn�es de $Vect1$ s�lectionn�es, ou $val1$ (m�me unit� que $Vect1$)\\
$Vect2_{sel}$ & Vecteur, \texttt{float} & Donn�es de $Vect2$ s�lectionn�es, ou $val2$ (m�me unit� que $Vect2$)\\
\end{tabarg}


%%%%%FORMULE%%%%
\itemstyle{Formule}

%%select_inter_data



%%%%%AUTEUR%%%%
\itemstyle{Commentaire~:}
Pour les valeurs manquantes du
vecteur �crit�re�, les valeurs des 2 vecteurs-r�sultats
sont manquantes.

%%%%%AUTEUR%%%%
\itemstyle{Auteur~:}
CNRM/GMEI/TRAMM

%%%%%REFERENCE%%%%
\itemstyle{R\'ef\'erences~:}

